\section{Continuity at ideal connections}\label{sec:continuity}

An ideal connection is a pair $(A_F,Z)$, where $F$ is a polystable locally
free sheaf with its Hermitian--Yang--Mills connection and
$Z=\sum_jn_jx_j$ is an effective zero-cycle. The determinant is fixed and
$c_2(F)+\deg Z=c_2(E)$.
We write $\MU$ for the Uhlenbeck compactification of the chosen moduli
component. Its algebraic structure is the reduced closure of that component
in Li's projective slope-semistable compactification, as used in
\cite[Section~7]{Taj23}. Only pairs in this closure are included.

\subsection{Algebraic and analytic convergence}
The comparison theorem needed here concerns the K\"ahler form used to solve
the Hermitian--Yang--Mills equation, not only its cohomology class.
We record the precise consequence.

\begin{lemma}\label{rec:comparison}
Let $[\omega]=c_1(H)$, with any smooth K\"ahler representative.
The algebraic and gauge-theoretic Uhlenbeck compactifications on a surface
are homeomorphic. If $(F_i,Z_i)\to(F,Z)$, then, after unitary gauge changes,
$A_{F_i}\to A_F$ smoothly on compact subsets away from $\Supp Z$, and
\begin{equation}\label{eq:energy-convergence}
 |F_{A_{F_i}}^0|^2dV_X+8\pi^2 Z_i
 \rightharpoonup |F_{A_F}^0|^2dV_X+8\pi^2 Z
\end{equation}
as measures. Here a zero-cycle denotes its atomic measure, with its
algebraic multiplicities.
\end{lemma}

\begin{proof}
In rank two, continuity for flat families for arbitrary Gauduchon metrics
is stated in \cite[Theorem~2.2]{BTT17}. For the arbitrary-rank surface
homeomorphism and convergence assertion we adapt
\cite[Theorem~4.11 and Corollary~4.15]{GSTW21} to the chosen representative.
We describe the metric-dependent steps of this adaptation.

Replace $H$ by a very ample power $L=H^a$, which does not change slope
stability. The $\partial\bar\partial$-lemma supplies a metric on $L$
with curvature $a\omega$. The bounded family of semistable sheaves admits
a common regularity twist by $L$. For every fixed sufficiently high twist,
choose a full $L^2$-orthonormal basis of sections. The Bochner inequality
for an HYM bundle twisted by $L^m$ gives
$\Delta|s|^2\le C_m|s|^2$ and therefore
$\|s\|_\infty\le C'_m\|s\|_2$. These estimates use the curvature
$a\omega$, not a Fubini--Study identity.
After Uhlenbeck convergence, every basis section has a subsequence
converging smoothly away from the bubbling set. The supremum estimate and
the negligible volume of small neighbourhoods of that set imply
\[
 \langle s_{\infty,\alpha},s_{\infty,\beta}\rangle_{L^2}
                    =\delta_{\alpha\beta}.
\]
Thus the whole basis survives, preserving the Hilbert-polynomial
inequality needed to identify the image sheaf. The stabilization and
full-rank arguments in \cite[Lemmas~4.3--4.4]{GSTW21} identify the Quot
limit; Proposition~4.5 there identifies its slope-graded double dual.
All slope comparisons involve the class $c_1(H)$ alone.

The individual multiplicities are compared by Chern--Weil transgression
on local slices in \cite[Propositions~4.7 and 4.9]{GSTW21}. These steps
use smooth K\"ahler geometry and local convergence, with no requirement
that the metric be induced by the projective embedding. Writing $E_0$ for the algebraic torsion-free limit, they identify
\[
 F=(\operatorname{gr}^{\mu}E_0)^{**},\qquad
 Z=\operatorname{cyc}\big(F/\operatorname{gr}^{\mu}E_0\big).
\]
The corrected definition in \cite[Erratum]{SW15} assigns the generic
associated-graded multiplicity, which on a surface is precisely the
length of this zero-dimensional quotient; it is also the convention in
\cite[Section~2.5.3]{GSTW21}.
Finally, the compactness and diagonal argument of
\cite[Proposition~4.13 and Theorem~4.11]{GSTW21} extends the comparison to
sequences of ideal connections and descends it through the quotient.
Corollary~4.15 gives the homeomorphism on surfaces. Uhlenbeck convergence
then gives the smooth convergence and measure identity in the statement.
\end{proof}

\subsection{Auxiliary data near a boundary point}
Fix $p=(F,Z)\in\MU$. A sufficiently high, general member $C\in|H^k|$
is smooth, avoids $\Supp Z$, and restricts each stable summand of $F$
to a stable bundle \cite{MR84}. Thus $F|_C$ is polystable, and
$\End F|_C$ is semistable of degree zero. Faltings' criterion \cite[Theorem~1.2]{Falt93}, in the form
\cite[Theorem~1.1]{BH08}, supplies a vector bundle $G$ on $C$ such that
\begin{equation}\label{eq:acyclic-limit}
 H^*(C,\End F|_C\otimes G)=0.
\end{equation}
Choose any smooth Hermitian metric on $G$ and a fixed metric on $C$.
These choices depend on $p$; there is no requirement that one curve or one
twist work on the whole compactification.

Choose a neighbourhood $N$ of $C$ with compact closure disjoint from
$\Supp Z$. Shrinking around $p$ in $\MU$, every nearby ideal cycle avoids
$\overline N$, and the restricted Dolbeault operator remains invertible.
Indeed, smooth convergence near $\overline N$ follows from
Lemma~\ref{rec:comparison}. If atoms occurred in $\overline N$ along a converging sequence, a
subsequence of their positions would converge there. Every such atom has
mass at least $8\pi^2$, forcing an atom of that mass in the limiting
measure, which is smooth on a neighbourhood of $\overline N$. This
contradiction excludes the atoms.
Invertibility is open in the smooth operator topology on the compact curve.
We may consequently use the same $C,G,N$ throughout an open neighbourhood
$U$ of $p$.

\begin{theorem}\label{thm:continuous-potential}
On $U$, define
\begin{align}\label{eq:ideal-potential}
 u(F',Z')={}&-\frac{\pi^2}{rk\rho}\log\det_\zeta\Delta_{\End F'|_C\otimes G}
       -\frac12\int_Xg_C|F_{A_{F'}}^0|^2dV_X\notag\\
 &\hspace{25mm}-4\pi^2\sum_jn'_jg_C(x'_j),
 \qquad Z'=\sum_jn'_jx'_j.
\end{align}
This function is continuous. Its restriction to every smooth holomorphic
family of stable bundles in $U$ is the potential of
Proposition~\ref{prop:potential}.
\end{theorem}
\begin{proof}
Let $(F_i,Z_i)\to(F',Z')$ in $U$. Choose a smooth cutoff $\chi$ supported in
$N$ and equal to one near $C$. On $\Supp\chi$, the curvatures converge
smoothly modulo gauge. Since $g_C\in L^1(X)$,
\begin{equation}\label{eq:near-curve}
 \int_X\chi g_C|F_{A_{F_i}}^0|^2dV_X
 \longrightarrow\int_X\chi g_C|F_{A_{F'}}^0|^2dV_X.
\end{equation}
The function $(1-\chi)g_C$ is continuous on all of $X$.
Applying \eqref{eq:energy-convergence} to this function gives
\begin{align*}
 &\int_X(1-\chi)g_C|F_{A_{F_i}}^0|^2dV_X
       +8\pi^2\sum_x\operatorname{mult}_x(Z_i)g_C(x)\\
 &\qquad\longrightarrow
 \int_X(1-\chi)g_C|F_{A_{F'}}^0|^2dV_X
       +8\pi^2\sum_x\operatorname{mult}_x(Z')g_C(x).
\end{align*}
Here $\chi$ vanishes at every atom under consideration.
Together with \eqref{eq:near-curve}, this proves continuity of the last two
terms in \eqref{eq:ideal-potential}.

On the fixed curve, the operators converge smoothly and remain invertible.
Their zeta determinants therefore converge
\cite[Theorem~1.6 and Remark~1.10]{BGS88}.
The determinant is gauge invariant, so the gauge changes used to obtain
smooth convergence do not change its value.
This proves continuity of $u$. On the locus of bundles the cycle is empty,
and the assertion about its complex Hessian is
Proposition~\ref{prop:potential}.
\end{proof}

The two parts of the proof address different features of bubbling.
The logarithmic singularity of $g_C$ is harmless where curvature converges
smoothly. Away from the curve the test function is continuous, so weak
convergence of the curvature measures retains exactly the atomic term in
\eqref{eq:ideal-potential}. No rate of convergence and no gluing expansion
are used.

The theorem gives a continuous function on the compactification. Positivity
of its complex Hessian on the normalization, and its global cohomology
class, require the arguments in Sections~\ref{sec:compactification}
and~\ref{sec:canonical}. At locally free polystable points it supplies a
surface version of the continuous-potential extension of
\cite[Theorem~11]{BSc20}; the additional term above treats nonempty bubbling
cycles.
